% !TEX root = main_AAAI.tex
\documentclass[letterpaper]{article}
\usepackage[submission]{aaai2027}
\usepackage[hyphens]{url}
\usepackage{graphicx}
\urlstyle{rm}
\def\UrlFont{\rm}
\usepackage{natbib}
\usepackage{caption}
\frenchspacing

\usepackage{amsmath,amssymb}
\usepackage{booktabs}

\pdfinfo{
/TemplateVersion (2027.1)
}
\setcounter{secnumdepth}{2}
\begin{document}
\title{COAST: Coordinated Ownership and Anticipatory Sequential Routing for Dynamic Vehicle Routing with Time Windows}
\author{Anonymous Submission}
\affiliations{}

	
\maketitle

\begin{abstract}
%Dynamic Vehicle Routing Problems with Time Windows (DVRPTW) require sequential policies to simultaneously coordinate fleet-wide decisions and anticipate long-term routing consequences under uncertainty. Existing shared policies often entangle these two aspects within a single scoring mechanism, leading to suboptimal decisions and limited generalization. We propose COAST, a neural routing framework that explicitly decomposes customer selection into two signals: a coordination-driven ownership bias derived from per-vehicle memory, and a candidate-conditioned lookahead estimate of future cost. These signals are integrated through an edge-aware scoring function that captures dynamic feasibility. Experiments on DVRPTW benchmarks show that COAST outperforms classical heuristics and recent neural baselines. Ablation studies demonstrate that coordination and anticipation provide complementary benefits, while out-of-distribution evaluations confirm improved generalization across varying arrival rates, time-window constraints, and fleet sizes. Our results highlight the importance of explicitly separating coordination and anticipation in sequential routing policies for dynamic environments.
Dynamic Vehicle Routing Problems with Time Windows (DVRPTW) require sequential routing policies to make decisions that are simultaneously feasible in the short term, coordinated across multiple vehicles, and effective over long planning horizons under dynamic customer arrivals. Existing neural routing policies typically compress these heterogeneous objectives into a single compatibility score, making coordination and anticipation difficult to learn jointly. We propose COAST (Coordinated Ownership and Anticipatory Sequential Routing), a neural routing framework that reformulates customer selection as a decision decomposition problem. Specifically, COAST estimates three complementary decision signals: Edge-aware Compatibility for immediate routing feasibility, Ownership Memory for implicit fleet coordination, and Candidate-conditioned Lookahead for anticipatory decision making. These signals are adaptively fused to produce the final routing policy and are optimized end-to-end through reinforcement learning without auxiliary supervision. Extensive experiments on standard DVRPTW benchmarks demonstrate that COAST consistently outperforms both classical heuristics and recent neural routing methods. Ablation studies confirm that the three decision signals provide complementary contributions, while out-of-distribution evaluations show substantially improved generalization across varying customer arrival rates, time-window tightness, and fleet sizes. These results suggest that explicitly decomposing routing decisions into complementary reasoning components provides a more effective and generalizable framework for dynamic vehicle routing.
\end{abstract}

\section{Introduction}\label{sec:intro}

The Vehicle Routing Problem (VRP) is one of the most extensively studied combinatorial optimization problems, with broad applications in logistics, transportation, and last-mile delivery. Classical VRP formulations assume that all customer requests and travel information are available before optimization begins \cite{toth2014vehicle}. However, many real-world logistics systems operate in highly dynamic environments, where customer requests arrive continuously during execution, traffic conditions evolve over time, and routing decisions must be updated online. These practical requirements give rise to the Dynamic Vehicle Routing Problem (DVRP), which more accurately models applications such as on-demand delivery, ride-sharing, and urban logistics \cite{psaraftis1988dynamic}.

Among its variants, the Dynamic Vehicle Routing Problem with Time Windows (DVRPTW) is particularly challenging because each request is associated with temporal and capacity constraints while future customer arrivals remain unknown. Unlike static routing, where globally optimized routes can be computed offline, a routing policy must repeatedly select the next customer for each vehicle under incomplete information. Decisions that appear locally optimal may adversely affect future routing opportunities, making effective online planning essential.

Traditional approaches to DVRPTW primarily rely on dispatching rules, metaheuristics with periodic re-optimization, or rolling-horizon optimization frameworks. Although these methods have demonstrated competitive performance in many settings, they often require extensive problem-specific engineering and repeated optimization during execution, making them difficult to generalize across different dynamic environments. Recently, deep reinforcement learning (DRL) has emerged as a promising alternative by learning routing policies directly from data, enabling real-time decision making without solving a new optimization problem whenever the environment changes.

Despite this progress, learning-based routing policies still struggle to simultaneously coordinate multiple vehicles and anticipate the long-term consequences of online decisions. Existing sequential routing policies typically employ a shared neural policy to make decisions for all vehicles. While this formulation naturally enables parameter sharing and implicit coordination, it also requires the policy to balance fleet-level coordination and long-term planning within a single decision process. As a result, the learned policy may become less robust when customer arrival patterns, spatial distributions, or time-window characteristics differ from those observed during training.

In this paper, we argue that coordination and long-term planning represent two fundamentally different reasoning processes and should be modeled explicitly rather than implicitly. Motivated by this observation, we propose COAST, a reinforcement learning framework that separates customer selection into complementary decision signals. COAST combines an edge-aware compatibility module for modeling local vehicle--customer interactions, a memory-based ownership mechanism for implicit fleet coordination, and a candidate-conditioned lookahead estimator for evaluating the downstream consequences of individual routing decisions. These complementary signals are adaptively fused into a unified routing policy that is trained end-to-end using reinforcement learning.

The main contributions of this paper are summarized as follows:

\begin{itemize}
    \item We present a structured decision decomposition framework for DVRPTW that explicitly separates fleet coordination from long-term planning within sequential routing policies.
    \item We introduce a memory-based ownership mechanism that enables implicit vehicle coordination without explicit inter-vehicle communication.
    \item We propose a candidate-conditioned lookahead module that estimates the long-term routing impact of each feasible action at the decision level.
    \item Extensive experiments demonstrate that the proposed framework consistently outperforms classical heuristics, hyper-heuristics, and recent neural routing approaches while exhibiting improved generalization across diverse dynamic environments.
\end{itemize}

The remainder of this paper is organized as follows. Section~\ref{sec:related} reviews the related literature. Section~\ref{sec:prob} formulates the problem. Section~\ref{sec:method} presents the proposed COAST framework. Section~\ref{sec:experiment} reports experimental results, and Section~\ref{sec:futurework} concludes the paper.



\section{Related Work}
\label{sec:related}

Deep Reinforcement Learning (DRL) has emerged as a competitive paradigm for neural combinatorial optimization in VRPs \cite{nazari2018reinforcement}. Early studies employed recurrent encoder--decoder architectures trained with policy gradients to learn routing heuristics directly from data \cite{nazari2018reinforcement}. Subsequently, attention-based models replaced recurrent networks with Transformer architectures, substantially improving representation learning and solution quality across a wide range of VRP variants \cite{kool2018attention}. More recent research has focused on enhancing the scalability, generalization, and expressive power of neural routing policies. Representative advances include transferable local policies for improving cross-scale generalization \cite{gao2023towards}, heterogeneous attention mechanisms for richer feature extraction \cite{wang2024deep}, lightweight cross-attention architectures for large-scale routing \cite{luo2025boosting}, and edge-aware Transformers that explicitly exploit pairwise customer relationships \cite{meng2025eformer}. These methods have significantly strengthened neural routing performance by improving representation learning and attention mechanisms. Nevertheless, routing decisions are still typically generated through a single compatibility function between vehicle and customer representations, implicitly coupling local feasibility, fleet coordination, and long-term routing considerations into a single decision score.

For DVRP, reinforcement learning has become an attractive alternative to repeatedly solving optimization models online. Existing studies formulate DVRPs as Markov decision processes and learn dispatching policies using value-based or actor--critic reinforcement learning \cite{akkerman2025comparison,konovalenko2024optimizing}. Recent work has further investigated richer state representations, graph neural networks, and Transformer-based architectures to better capture dynamic customer arrivals and evolving routing states \cite{konovalenko2024optimizing,li2024cada,meng2025eformer}. Although these advances improve policy quality and generalization, routing decisions are still primarily driven by a unified compatibility score, making it difficult to explicitly distinguish immediate feasibility from the long-term consequences of selecting a particular customer.

Another fundamental challenge is coordinating multiple vehicles in dynamic environments. Existing approaches typically rely on centralized training, explicit inter-vehicle communication, or multi-agent reinforcement learning to coordinate vehicle behaviors \cite{zhang2020multi,shahbazian2025knowledge,xia2025multi}. While these strategies improve cooperation, they generally require repeatedly modeling the entire fleet during decision making, leading to increased computational complexity and reduced scalability. Furthermore, most actor--critic methods estimate only a state-value function for variance reduction. Since the critic is independent of the selected action, it provides identical guidance for all candidate customers within the same state, limiting its ability to evaluate the downstream routing impact of individual actions and making learned policies more susceptible to myopic decisions.

Taken together, existing learning-based routing policies exhibit three important limitations: (i) local feasibility, fleet coordination, and long-term planning are entangled within a single compatibility score; (ii) effective vehicle coordination often relies on computationally expensive fleet-level interaction; and (iii) state-value critics provide action-independent guidance that cannot explicitly assess the long-term impact of individual routing decisions. To address these limitations, we propose COAST, which decomposes customer selection into three complementary signals: an edge-aware compatibility module for local feasibility, a memory-based ownership mechanism for implicit fleet coordination, and a candidate-conditioned lookahead estimator for evaluating downstream routing consequences. These signals are adaptively fused into a unified routing policy trained end-to-end using reinforcement learning.
\section{Problem Statement}
\label{sec:prob}

\subsection{Problem Setting}
\label{sec:prob_setting}

A DVRPTW instance consists of a depot (node $0$) and a set of customer requests
$\mathcal{C}=\{1,\ldots,L_c\}$. Each customer request $j$ is represented by
\[
\mathbf{x}_j=(x_j,y_j,q_j,e_j,l_j,s_j),
\]
where $(x_j,y_j)$ denotes the normalized location, $q_j$ is the normalized demand,
$[e_j,l_j]$ is the service time window, and $s_j$ is the service duration.

At decision step $t$, vehicle $k$ is described by
\[
\mathbf{s}_k^t=(x_k,y_k,q_k^{rem},t_k^{cur}),
\]
where $(x_k,y_k)$ is the current position, $q_k^{rem}$ is the remaining capacity,
and $t_k^{cur}$ is the next available time.

For assigning vehicle $k$ to customer $j$, the travel distance, travel time,
arrival time, and tardiness are computed as
\begin{align}
d_{kj} &= \|(x_k,y_k)-(x_j,y_j)\|_2,\\
\tau_{kj} &= d_{kj}/v,\\
t^{arr}_{kj} &= t_k^{cur}+\tau_{kj},\\
\ell_{kj} &= \max(0,t^{arr}_{kj}-l_j).
\end{align}

Following~\cite{TAM2026114234}, we adopt soft time windows, allowing late arrivals
while penalizing tardiness in the objective. A customer is feasible whenever
$q_j\le q_k^{rem}$.

Unlike the static VRPTW, customer requests are revealed progressively during route
execution. Initially only $\mathcal{C}_0$ is observable, while additional batches
$\mathcal{C}_{e_1},\mathcal{C}_{e_2},\ldots$ arrive over time. Let
\[
\mathcal{C}_t=\bigcup_{e_i\le t}\mathcal{C}_{e_i}
\]
denote the set of currently revealed requests. The degree of dynamism (DoD)
measures the proportion of dynamically arriving customers. Consequently, routing
decisions must balance serving current requests while preserving sufficient
capacity and time for future arrivals.

\subsection{Sequential Multi-Agent Decision Process}
\label{sec:prob_mdp}

Following~\cite{TAM2026114234}, we formulate DVRPTW as a Sequential Multi-Agent
Decision Process (sMDDP), defined by
\[
(\mathcal I,\mathcal S,\mathcal A,P,P_0,R,T),
\]
where $\mathcal I=\{1,\ldots,L_v\}$ is the homogeneous vehicle set.

Instead of selecting joint actions for all vehicles simultaneously, only the
vehicle that becomes available first makes a decision at each step. Let
$\iota_t\in\mathcal I$ denote this active vehicle.

The global state is
\[
s_t=(\{\mathbf{s}_k^t\}_{k\in\mathcal I},\,s_t^E,\,\iota_t),
\]
where $s_t^E$ stores the status of all revealed customers.

At decision step $t$, the action is
\[
a_t=j,\qquad
j\in\mathcal A(s_t,\iota_t),
\]
where $\mathcal A(s_t,\iota_t)$ contains all revealed and unserved customers
satisfying the capacity constraint together with the depot action.

The transition function
\[
P(s_{t+1}\mid s_t,a_t)
\]
updates the active vehicle state, remaining capacity, customer status, newly
revealed requests, and selects the next active vehicle according to the earliest
available time.

All vehicles share a common policy
\[
\pi_\theta(a_t\mid s_t,\iota_t),
\]
allowing routing knowledge to generalize across the fleet while naturally
supporting different fleet sizes during inference.

\subsection{Objective}
\label{sec:prob_obj}

The routing objective jointly minimizes travel distance, service tardiness, and
the number of unserved customers.

For serving customer $a_t$, the immediate reward is
\[
r_t
=
-d_{\iota_t,a_t}
-
\lambda_{late}\ell_{\iota_t,a_t},
\]
while the terminal penalty is
\[
r_T^{skip}
=
-\lambda_{skip}N_{skip},
\]
where $N_{skip}$ denotes the number of revealed but unserved customers.

The cumulative return is
\[
\begin{aligned}
R
&= \sum_{t=0}^{T-1}r_t+r_T^{skip} \\
&= -\sum_t d_{\iota_t,a_t}
   -\lambda_{late}\sum_t\ell_{\iota_t,a_t}
   -\lambda_{skip}N_{skip}.
\end{aligned}
\]

The training objective is
\[
\max_\theta \mathbb{E}_{\pi_\theta}[R].
\]

This formulation highlights three key challenges: (i) dynamic customer arrivals
lead to a non-stationary decision process; (ii) routing decisions of different
vehicles are sequentially coupled through shared customers; and (iii) routing
feasibility depends on the interaction between vehicle states and customer
attributes. These challenges directly motivate the edge-aware compatibility,
ownership coordination, and candidate-conditioned lookahead mechanisms proposed
in the next section.

\section{Methodology}\label{sec:method}

This section presents the COAST architecture and its training procedure.
Table~\ref{tab:notation} summarizes the main notation.

\subsection{Policy Overview}\label{sec:policy_overview}

Existing neural routing policies typically produce a selection score through a single compatibility function between vehicle and customer representations~\cite{kool2019attention}, which conflates immediate feasibility assessment with fleet-level coordination and future planning.
When these signals interfere, the policy gradient receives mixed learning signals that slow convergence and limit generalization under distribution shifts.

COAST instead decomposes the customer-selection score into three independently computed signals combined by a learned fusion mechanism.
At each step $t$, given state $s_t$ and active vehicle $k_t$, the policy is:
\[
\pi_\theta(a_t \mid s_t, k_t)
=
\frac{\exp\!\big(S_\theta(a_t \mid s_t, k_t)\big)}
{\displaystyle\sum_{u \in \mathcal{A}(s_t, k_t)} \exp\!\big(S_\theta(u \mid s_t, k_t)\big)},
\]
where $\mathcal{A}(s_t, k_t)$ is the feasible action set and the composite score is:
\[
S_\theta(j \mid s_t, k_t) = \phi_{fuse}\!\left(\tilde{S}^{att}_j,\; \tilde{S}^{owner}_j,\; \tilde{S}^{look}_j\right).
\]
The three signals are:
\begin{itemize}
    \item $S^{att}$: a \textit{local compatibility score} quantifying vehicle--customer fit given current spatial, temporal, and capacity constraints;
    \item $S^{owner}$: a \textit{coordination bias} derived from per-vehicle memory, encoding implicit fleet-level assignment of customer $j$ to vehicle $k_t$;
    \item $S^{look}$: a \textit{candidate-conditioned lookahead estimate} of the downstream routing consequence of committing to customer $j$.
\end{itemize}

\subsection{Model Architecture}\label{sec:model_architecture}

\paragraph{Customer graph encoding.}
We encode customers using a Transformer-based graph encoder~\cite{vaswani2017attention} that produces contextual embeddings capturing inter-customer relational structure.
Depot and customer nodes use separate linear projections since they play structurally different roles:
\[
\mathbf{h}_j^{(0)} =
\begin{cases}
\mathbf{W}_{dep}\,\mathbf{x}_j, & j = 0 \\
\mathbf{W}_{cust}\,\mathbf{x}_j, & j \geq 1
\end{cases}
\]
where $\mathbf{x}_j \in \mathbb{R}^{D_c}$ is the raw feature vector.
Since geographic proximity is a strong signal for routing quality but may not emerge reliably from content-based attention alone, we augment each attention layer with a distance-aware bias: a 16-bin RBF encoding of normalized pairwise distance $d_{ij} \in [0,1]$,
\[
\mathrm{RBF}_k(d) = \exp\!\left(-\frac{(d - \mu_k)^2}{2\sigma^2}\right),
\]
with 16 evenly-spaced centres $\mu_k = (k-1)/15$, $k=1,\ldots,16$, and bandwidth $\sigma = 1/15$.
The 16-dimensional encoding is projected by an MLP to a per-head scalar bias $b_{ij}^{(l,h)}$ added to the attention logit.
Each layer applies two post-norm residual sublayers~\cite{vaswani2017attention}.
Letting $\mathrm{MHA}_i^{(l)} = \mathbf{W}_{out}^{(l)}\!\left[\bigoplus_{h=1}^{H}\sum_j \alpha_{ij}^{(l,h)}\mathbf{v}_j^{(l,h)}\right]$ denote the multi-head attention output:
\begin{align*}
\hat{\mathbf{h}}_i^{(l)} &= \mathrm{LN}\!\left(\mathbf{h}_i^{(l)} + \mathrm{MHA}_i^{(l)}\right), \\
\mathbf{h}_i^{(l+1)} &= \mathrm{LN}\!\left(\hat{\mathbf{h}}_i^{(l)} + \mathrm{FFN}\!\left(\hat{\mathbf{h}}_i^{(l)}\right)\right).
\end{align*}
The final customer representation is $\mathbf{c}_j = \mathrm{Dropout}(\mathbf{W}_c\,\mathbf{h}_j^{(L)}) \in \mathbb{R}^D$.
In the dynamic setting, encodings are recomputed whenever new orders become visible.

\paragraph{Acting vehicle encoding.}
Only the active vehicle $k_t$ is encoded at each step --- full fleet re-encoding at every step is wasteful and unnecessary since inter-vehicle coordination is handled by CoordinationMemory rather than vehicle-to-vehicle attention.
The vehicle state $\mathbf{s}_{k_t} = (x, y, q^{rem}, t^{cur}) \in \mathbb{R}^4$ is projected and refined through $L$ layers of cross-attention over $H^{cust}$, restricted to the feasible action set:
\[
h^{act} = \mathbf{W}_v\;\mathrm{FleetEncoder}\!\bigl(\mathbf{W}_{in}\,\mathbf{s}_{k_t},\; H^{cust},\; \mathcal{A}(s_t,k_t)\bigr),
\]
where $h^{act} \in \mathbb{R}^{N \times 1 \times D}$.

\paragraph{Edge feature encoding}
Node-level representations capture global structure but cannot express the state-dependent interaction between a specific vehicle and a specific candidate at this step.
Two customers equidistant from the vehicle may differ radically in time-window compatibility, waiting time, or remaining capacity fit.
We therefore construct an 8-dimensional edge feature vector $\mathbf{e}_{ij} \in \mathbb{R}^8$ for each (vehicle $k_t$, candidate $j$) pair:

\medskip
\begin{center}
\begin{tabular}{@{} l l p{1.65in} @{}}
\toprule
Feature & Symbol & Definition \\
\midrule
Distance       & $d_{ij}$    & $\|\mathbf{p}_{k_t} - \mathbf{p}_j\|_2$ \\
Travel time    & $\tau_{ij}$ & $d_{ij} / v$ \\
Arrival time   & $a_{ij}$    & $t^{cur}_{k_t} + \tau_{ij}$ \\
Waiting time   & $w_{ij}$    & $\max(tw^{start}_j - a_{ij},\,0)$ \\
Lateness       & $l_{ij}$    & $\max(a_{ij} - tw^{end}_j,\,0)$ \\
Time slack     & $s_{ij}$    & $tw^{end}_j - a_{ij}$ \\
Feasibility    & $f_{ij}$    & 1 if $j$ is capacity- and time-feasible, else 0 \\
Capacity gap   & $g_{ij}$    & $q^{rem}_{k_t} - d(j)$ \\
\bottomrule
\end{tabular}
\end{center}
\medskip

Together these features encode whether the vehicle arrives on time, how much schedule slack remains, and whether capacity is sufficient --- information that changes at every step as the vehicle's state evolves.
The features are encoded through a two-layer MLP with layer normalization:
\[
\mathbf{e}_{ij}^{emb} = \mathrm{LN}\!\left(\mathbf{W}_{e2}\,\mathrm{ReLU}\!\left(\mathbf{W}_{e1}\,\mathbf{e}_{ij}\right)\right) \in \mathbb{R}^D,
\]
yielding the full edge tensor $H^{edge} \in \mathbb{R}^{N \times 1 \times L_c \times D}$.

\paragraph{Local compatibility via edge-aware scoring.}
Standard attention scoring~\cite{kool2019attention} computes vehicle--customer compatibility from node-level representations that are oblivious to the specific interaction at this step.
The CrossEdgeFusion module injects the per-pair edge embedding as an additive bias into multi-head attention logits:
\[
u_j^{(h)} = \frac{\mathbf{q}^{(h)} \cdot \mathbf{k}_j^{(h)}}{\sqrt{d_h}} + \mathbf{e}_j^{emb} \cdot \mathbf{w}_{bias}^{(h)},
\]
where $\mathbf{q}^{(h)} = h^{act}\mathbf{W}_Q^{(h)}$ and $\mathbf{k}_j^{(h)} = \mathbf{c}_j\mathbf{W}_K^{(h)}$.
The compatibility score averages the pre-softmax logits across heads:
\[
S^{att}_j = \frac{1}{H}\sum_{h=1}^{H} u_j^{(h)} \in \mathbb{R}^{N \times 1 \times L_c}.
\]
No softmax is applied inside CrossEdgeFusion; the single softmax over all candidates is deferred to the final action selection stage after all three signals are fused, ensuring that inter-candidate probability mass is distributed over the most informative composite signal.

\paragraph{Coordination through memory and ownership.}
Without a coordination mechanism, the policy may route multiple vehicles toward the same customer clusters while neglecting others.
Explicit fleet-level communication at every step is computationally heavy and unnecessary when softer implicit signals suffice.

We instead equip each vehicle with a private latent memory $\mathbf{m}_k^{(t)} \in \mathbb{R}^{H_m}$, initialized to zero, that summarizes its routing history.
After vehicle $k_t$ selects customer $j^*$, its memory is updated via a bounded, GRU-like recurrence:
\[
\mathbf{m}_{k_t}^{(t+1)} = \tanh\!\left(\mathbf{W}_{in}\!\left[\,h^{act},\;\mathbf{c}_{j^*},\;\mathbf{e}_{k_t j^*}^{emb}\,\right] + \mathbf{W}_{hid}\,\mathbf{m}_{k_t}^{(t)}\right).
\]
Only the acting vehicle's memory slot is updated; all others remain unchanged.

This per-vehicle memory supports an OwnershipHead that computes a soft fleet-level customer assignment.
A vehicle whose memory is strongly aligned with customer $j$ has implicitly claimed it through prior routing choices.
The assignment logits are computed via a bilinear interaction over all vehicles:
\[
o_{ik} = \frac{(\mathbf{W}_{veh}\,\mathbf{m}_i^{(t)}) \cdot (\mathbf{W}_{cust}\,\mathbf{c}_k)^\top}{\sqrt{D}},
\]
and the coordination bias for vehicle $k_t$ is the log-relative ownership over the fleet:
\[
\begin{aligned}
S^{owner}_j
&= \log\!\left(\frac{\exp(o_{k_t,j})}{\sum_{i'=1}^{L_v}\exp(o_{i',j})}\right), \\
S^{owner} &\in \mathbb{R}^{N \times 1 \times L_c}.
\end{aligned}
\]
Using log-probability ensures the ownership signal is well-scaled relative to the other scores, and the softmax across vehicles encodes \textit{relative} claim --- a high $S^{owner}_j$ means vehicle $k_t$ is the strongest claimant among all vehicles.
Neither the OwnershipHead nor the memory update imposes any auxiliary supervised loss; both are shaped entirely by the policy gradient, allowing coordination semantics to emerge from optimizing routing performance.

\paragraph{Candidate-conditioned lookahead estimation.}
A policy that maximizes immediate feasibility may still be myopic: selecting a nearby customer can leave the vehicle in a position that violates time windows of unserved neighbors, or forces an early depot return due to capacity exhaustion.
A global state value $V(s_t)$ cannot differentiate between candidates and therefore offers no guidance at the action level.

The LookaheadHead produces a per-candidate scalar that estimates the downstream routing consequence of committing to customer $j$:
\[
S^{look}_j = \mathbf{W}_{\ell 2}\;\mathrm{Dropout}\!\bigl(\mathrm{ReLU}\bigl(\mathbf{W}_{\ell 1}
\bigl[\,h^{act},\;\mathbf{c}_j,\;\mathbf{e}_{ij}^{emb}\,\bigr]\bigr)\bigr),
\]
producing a per-candidate scalar ($S^{look}_j \in \mathbb{R}$) conditioned on the vehicle's current state, the candidate's global embedding, and their dynamic interaction.
Critically, this head carries no explicit regression target --- $S^{look}$ is shaped entirely through the policy gradient, which drives the fusion MLP to weight lookahead appropriately when it reduces long-term routing cost.

\paragraph{Multi-signal score fusion.}
The three signals are computed by different mechanisms and operate on incompatible numerical scales.
To normalize them before fusion, we apply mask-aware z-normalization over the feasible candidates for each signal independently.
For each $j \in \mathcal{A}(s_t, k_t)$:
\[
\tilde{S}^{(\cdot)}_j = \frac{S^{(\cdot)}_j - \mu^{(\cdot)}}{\sqrt{(\sigma^{(\cdot)})^2 + \varepsilon}},
\]
where $\mu^{(\cdot)}$ and $\sigma^{(\cdot)}$ are the mean and standard deviation computed \emph{only} over feasible candidates; scores for infeasible positions are subsequently set to $-\infty$ before the final softmax.

The normalized signals are fused by a compact two-layer MLP:
\[
S^{final}_j = \mathbf{W}_{f2}\;\mathrm{ReLU}\!\left(\mathbf{W}_{f1}\!\left[\,\tilde{S}^{att}_j,\;\tilde{S}^{owner}_j,\;\tilde{S}^{look}_j\,\right]\right),
\]
with $\mathbf{W}_{f1} \in \mathbb{R}^{64 \times 3}$ and $\mathbf{W}_{f2} \in \mathbb{R}^{1 \times 64}$.
A fixed-weight sum cannot adapt to the fact that ownership is uninformative early in an episode but becomes the dominant coordination cue later; the MLP learns these context-dependent interaction weights through the policy gradient.
An optional tanh scaling bounds the magnitude: $S^{final}_j \leftarrow \beta_0 \cdot \tanh(S^{final}_j)$.

\subsection{Training Objective}\label{sec:training_objective}

COAST is trained end-to-end using REINFORCE~\cite{nazari2018reinforcement} with a learned critic baseline.
The critic $b_\phi$ and the LookaheadHead play strictly separate roles: $b_\phi$ is a state-level variance-reducing baseline that is agnostic to the action chosen, while $S^{look}$ is an action-conditioned routing guide.
Conflating them would couple value estimation with the scoring mechanism, degrading both.

\paragraph{Policy gradient with critic baseline.}
The policy loss is:
\[
\mathcal{L}_{policy}(\theta)
=
-\mathbb{E}_{\tau \sim \pi_\theta}\!\left[\sum_{t=0}^{T} \log \pi_\theta(a_t \mid s_t, k_t)\, A_t \right],
\]
where $A_t = R_t - b_\phi(s_t)$ is the advantage, $R_t = \sum_{k=t}^T r_k$ is the cumulative return, and the critic is trained with:
\[
\mathcal{L}_{critic}(\phi)
=
\mathbb{E}\!\left[\sum_{t=0}^{T} \mathrm{SmoothL1}\!\left(b_\phi(s_t),\; R_t\right)\right].
\]
Smooth-L1 is used for robustness to the return variance induced by stochastic customer arrivals.
Advantage estimates are batch-normalized before computing $\mathcal{L}_{policy}$.

\paragraph{Combined objective.}
We augment with an entropy term to prevent premature convergence across the diverse instance distribution:
\begin{align*}
\mathcal{L}(\theta, \phi)
&= \mathcal{L}_{policy}(\theta)
+ \alpha\; \mathcal{L}_{critic}(\phi) \\
&\quad -\,\beta\;\mathbb{E}\!\left[\sum_{t=0}^{T}
  \mathcal{H}\!\left(\pi_\theta(\cdot \mid s_t, k_t)\right)\right],
\end{align*}
with $\beta = 0.01$.
The LookaheadHead and OwnershipHead carry no auxiliary losses; both receive gradient signal exclusively through $S^{final}$ via the policy gradient.

\subsection{Discussion}\label{sec:method_discussion}

The design of COAST reflects three structural inductive biases aligned with dynamic multi-vehicle routing.

\textit{Edge features as first-class routing context.}
Static node embeddings encode global relational structure but cannot express how vehicle state and customer constraints interact at a given step.
The explicit construction of per-pair edge features --- encoding time slack, lateness, capacity gap, and feasibility --- allows the scoring mechanism to react immediately to changes in vehicle state without waiting for global re-encoding.
This is the primary mechanism through which COAST achieves constraint-awareness at decision time.

\textit{Communication-free fleet coordination via emergent ownership.}
Explicit inter-vehicle communication or full fleet re-encoding at every step adds overhead proportional to fleet size and introduces gradient coupling between vehicles.
The per-vehicle memory and OwnershipHead provide a scalable alternative: coordination emerges implicitly from each vehicle's trajectory history, naturally inducing spatial specialization without any vehicle-to-vehicle attention.
The ownership signal is shaped purely by routing performance, not by any auxiliary supervision, so the coordination structure it encodes reflects what actually improves fleet efficiency.

\textit{Action-conditioned lookahead without a supervised target.}
The LookaheadHead differs from standard value-based approaches in that it produces per-candidate estimates without being constrained to approximate any specific quantity such as $Q(s_t, a_t)$.
By training entirely through the policy gradient, the lookahead signal is free to encode whatever combination of vehicle state, candidate embedding, and edge interaction best reduces myopic decisions --- including routing consequences that are difficult to formalize analytically.
Its separation from the critic baseline preserves the action-independence required for an unbiased gradient estimator while allowing the per-candidate signal to carry independent predictive content.

\begin{table*}[t!]
\centering
{\small
\setlength{\tabcolsep}{2pt}
\begin{tabular}{p{0.57\textwidth}p{0.39\textwidth}}
\toprule
\textbf{Symbol} & \textbf{Description} \\
\midrule

\multicolumn{2}{l}{\textit{Decision process}} \\
$t$                          & Decision step index \\
$s_t$                        & Environment state at step $t$ \\
$k_t$                        & Index of the active vehicle at step $t$ \\
$a_t$                        & Selected customer (action) \\
$\mathcal{A}(s_t, k_t)$      & Feasible action set for vehicle $k_t$ \\

\midrule
\multicolumn{2}{l}{\textit{Problem dimensions}} \\
$N$   & Batch size \\
$L_c$ & Number of customers including depot \\
$L_v$ & Number of vehicles \\
$D$   & Embedding dimension \\
$H$   & Number of attention heads \\
$H_m$ & Coordination memory hidden size \\
$H_\ell$ & Lookahead MLP hidden size \\

\midrule
\multicolumn{2}{l}{\textit{Input features}} \\
$\mathbf{x}_j \in \mathbb{R}^{D_c}$      & Raw feature vector of customer or depot $j$ \\
$\mathbf{s}_{k_t} \in \mathbb{R}^4$       & State of active vehicle: $(x, y, q^{rem}, t^{cur})$ \\
$\mathbf{e}_{ij} \in \mathbb{R}^8$        & Raw edge feature vector for pair (vehicle $k_t$, candidate $j$) \\

\midrule
\multicolumn{2}{l}{\textit{Learned representations}} \\
$\mathbf{c}_j \in \mathbb{R}^D$,\; $H^{cust} \in \mathbb{R}^{N \times L_c \times D}$               & Customer embeddings from GraphEncoder \\
$h^{act} \in \mathbb{R}^{N \times 1 \times D}$                                                       & Active vehicle embedding (one vehicle encoded per step) \\
$\mathbf{e}_{ij}^{emb} \in \mathbb{R}^D$,\; $H^{edge} \in \mathbb{R}^{N \times 1 \times L_c \times D}$ & Encoded edge features \\

\midrule
\multicolumn{2}{l}{\textit{Scoring signals}} \\
$S^{att}   \in \mathbb{R}^{N \times 1 \times L_c}$ & Local compatibility score (CrossEdgeFusion) \\
$S^{owner} \in \mathbb{R}^{N \times 1 \times L_c}$ & Ownership coordination bias (log-probability) \\
$S^{look}  \in \mathbb{R}^{N \times 1 \times L_c}$ & Candidate-conditioned lookahead estimate \\
$S^{final} \in \mathbb{R}^{N \times 1 \times L_c}$ & Fused routing score \\
$\tilde{S}^{(\cdot)}$                               & Z-normalized version of each signal \\

\midrule
\multicolumn{2}{l}{\textit{Coordination memory}} \\
$\mathbf{m}_k^{(t)} \in \mathbb{R}^{H_m}$    & Memory state of vehicle $k$ at step $t$ \\
$M_t \in \mathbb{R}^{N \times L_v \times H_m}$ & Full fleet memory tensor at step $t$ \\

\midrule
\multicolumn{2}{l}{\textit{Policy and learning}} \\
$\pi_\theta(a_t \mid s_t, k_t)$  & Routing policy parameterized by $\theta$ \\
$R_t$                            & Cumulative return from step $t$ \\
$b_\phi(s_t) \in \mathbb{R}$     & Learned critic baseline (distinct from LookaheadHead) \\
$A_t = R_t - b_\phi(s_t)$        & Advantage estimate \\
$\alpha,\,\beta$                 & Critic weight and entropy coefficient \\
\bottomrule
\end{tabular}
}
\caption{Summary of main notation used in Section~\ref{sec:method}.}
\label{tab:notation}
\end{table*}

\section{Experimental results} \label{sec:experiment}

\subsection{Evaluation protocol}

We evaluate the normalized routing objective from Section~\ref{sec:prob}, which combines travel distance, lateness, and unserved-request penalties; lower is better. All neural methods use greedy decoding on fixed, matched test instances. The Solomon-style suite contains 176 instances across \texttt{h100}, \texttt{h200}, and \texttt{h400}; PYTH contains 10,240 instances at each scale. COAST and every internal variant use the same instance-generation and optimization protocol, with an ablation changing only its named design. MARDAM, AM, and PolyNet are external neural baselines; GPHH, C+C, C+W, and WIQ+C are hyper-heuristic or dispatching-rule baselines.

RQ1 keeps benchmark reporting and statistical inference distinct. On the Solomon-style suite, all eight methods are evaluated on the same 176 instances. We first apply a Friedman omnibus test across COAST, the three neural baselines, GPHH, and the three dispatching rules, then test the seven prespecified COAST--baseline contrasts with two-sided paired Wilcoxon tests and Holm correction. A neural-only Friedman test provides a like-for-like complement. These tests quantify consistency across matched test instances, not robustness across independently trained seeds.

\subsection{RQ1: Does COAST improve benchmark routing quality?}

\paragraph{Question.} Does the decision decomposition improve routing quality across spatial topologies, dispatching baselines, and increasing fleet scale?

Tables~\ref{tab:rq1_solomon_neural} and \ref{tab:rq1_solomon_classical} split the complete Solomon-style benchmark into learned and classical comparators. This preserves all R, C, and RC topology evidence while keeping each comparison set readable in a single column.

\begin{table}[t]
\centering
{\small
\setlength{\tabcolsep}{2.4pt}
\begin{tabular}{llrrrr}
\toprule
Scale & Top. & COAST & MARDAM & AM & Poly. \\
\midrule
\texttt{h100} & C  & \textbf{67.65} & 65.08 & 73.01 & 70.36 \\
               & R  & \textbf{57.22} & 60.03 & 66.08 & 63.97 \\
               & RC & \textbf{65.69} & 67.42 & 73.07 & 70.31 \\
\addlinespace
\texttt{h200} & C  & \textbf{117.61} & 162.43 & 126.64 & 137.04 \\
               & R  & \textbf{104.24} & 142.69 & 124.00 & 125.65 \\
               & RC & \textbf{105.19} & 143.48 & 125.66 & 130.48 \\
\addlinespace
\texttt{h400} & C  & \textbf{220.65} & 323.84 & 248.35 & 258.78 \\
               & R  & \textbf{201.98} & 314.44 & 243.97 & 238.60 \\
               & RC & \textbf{203.26} & 322.73 & 250.32 & 244.98 \\
\bottomrule
\end{tabular}
}
\caption{Solomon-style normalized routing cost: neural comparisons (lower is better). ``Poly.'' denotes PolyNet. COAST is best in 26/27 size--topology aggregates; MARDAM is lower only on \texttt{h100}-C.}
\label{tab:rq1_solomon_neural}
\end{table}

\begin{table}[t]
\centering
{\small
\setlength{\tabcolsep}{2.0pt}
\begin{tabular}{llrrrrr}
\toprule
Scale & Top. & COAST & GPHH & C+C & C+W & WIQ+C \\
\midrule
\texttt{h100} & C  & \textbf{67.65} & 80.63 & 130.71 & 118.12 & 175.02 \\
               & R  & \textbf{57.22} & 72.02 & 147.70 & 143.05 & 179.86 \\
               & RC & \textbf{65.69} & 80.11 & 152.99 & 141.10 & 181.66 \\
\addlinespace
\texttt{h200} & C  & \textbf{117.61} & 191.74 & 269.58 & 248.67 & 361.24 \\
               & R  & \textbf{104.24} & 184.87 & 276.62 & 273.03 & 348.95 \\
               & RC & \textbf{105.19} & 181.00 & 282.86 & 283.12 & 350.99 \\
\addlinespace
\texttt{h400} & C  & \textbf{220.65} & 374.85 & 630.06 & 616.13 & 752.28 \\
               & R  & \textbf{201.98} & 354.62 & 624.08 & 618.70 & 720.14 \\
               & RC & \textbf{203.26} & 352.72 & 634.84 & 636.76 & 724.18 \\
\bottomrule
\end{tabular}
}
\caption{Solomon-style normalized routing cost: GPHH and dispatching comparisons. COAST is lower-cost than GPHH on 172/176 instances and than each dispatching rule on all 176.}
\label{tab:rq1_solomon_classical}
\end{table}

\paragraph{Solomon benchmark.} The two Solomon tables establish three points. First, COAST is best in every \texttt{h200} and \texttt{h400} size--topology aggregate. Second, the small-scale exception is explicit rather than averaged away: at \texttt{h100}, MARDAM is 3.8\% lower on the compact C topology, while COAST still leads the \texttt{h100} aggregate and its R and RC subsets. Measured from the scale-level means, COAST's reduction over the best neural comparator is 1.9\%, 13.0\%, and 15.1\% at \texttt{h100}, \texttt{h200}, and \texttt{h400}, respectively.

Third, Table~\ref{tab:rq1_solomon_classical} places the neural comparison in a broader dispatching context. COAST beats GPHH on 172/176 individual instances and every dispatching rule on all 176. The separation is especially pronounced at \texttt{h400}: on R instances, COAST's 201.98 cost is below GPHH's 354.62 and the two savings-based rules' 624.08 and 618.70. Thus, the result is not limited to one learned comparator, while the neural table remains the more demanding comparison because its baselines are closer.

\begin{table*}[t]
\centering
\begin{minipage}[t]{0.47\textwidth}
\centering
{\small
\setlength{\tabcolsep}{3.0pt}
\begin{tabular}{lrccc}
\toprule
Topology & $n$ & \shortstack{MARDAM\\win / $\tilde{\Delta}$} & \shortstack{AM\\win / $\tilde{\Delta}$} & \shortstack{PolyNet\\win / $\tilde{\Delta}$} \\
\midrule
R  & 63 & 95.2\% / 23.1\% & 98.4\% / 15.4\% & 93.7\% / 14.8\% \\
C  & 57 & 82.5\% / 21.3\% & 89.5\% / 9.6\%  & 87.7\% / 11.0\% \\
RC & 56 & 89.3\% / 25.0\% & 94.6\% / 15.8\% & 92.9\% / 15.0\% \\
\bottomrule
\end{tabular}
}\\[1pt]
\footnotesize (a) Topology-stratified neural comparisons.
\end{minipage}\hfill
\begin{minipage}[t]{0.47\textwidth}
\centering
{\small
\setlength{\tabcolsep}{4.5pt}
\begin{tabular}{lrrrr}
\toprule
PYTH scale & COAST & MARDAM & AM & PolyNet \\
\midrule
\texttt{n20m1}$^{\dagger}$ & \textbf{24.78} & 26.13 & 25.40 & 25.58 \\
\texttt{n50m3}             & \textbf{43.92} & 46.48 & 49.58 & 49.93 \\
\texttt{n100m5}            & \textbf{94.85} & 99.50 & 103.85 & 105.13 \\
\texttt{n200m10}           & \textbf{182.38} & 203.24 & 196.53 & 200.89 \\
\texttt{n400m20}           & \textbf{355.99} & 430.48 & 378.98 & 389.72 \\
\bottomrule
\end{tabular}
}\\[1pt]
\footnotesize (b) PYTH greedy-decoding results.
\end{minipage}
\caption{Detailed RQ1 benchmark evidence. In (a), ``win / $\tilde{\Delta}$'' gives COAST's instance win rate and median paired percentage cost reduction against each neural comparator. In (b), each PYTH scale has 10,240 matched instances and reports mean normalized cost under greedy decoding; $^{\dagger}$\texttt{n20m1} is a single-vehicle control.}
\label{tab:rq1_detail}
\end{table*}

Table~\ref{tab:rq1_detail}(a) shows where the Solomon gains arise. C is consistently the closest topology: it has the lowest COAST win rate and median reduction for each of the three neural comparisons. Even there, COAST wins 82.5\% of matched instances against MARDAM and 87.7\% against PolyNet. The R and RC results are stronger, with 89.3--98.4\% win rates across the three neural comparators. The spatial evidence therefore does not depend on a single topology.

\paragraph{PYTH benchmark.} Table~\ref{tab:rq1_detail}(b) provides complementary scale evidence by holding the generated family fixed while jointly increasing customer count and fleet size; every row contains 10,240 matched instances. The single-vehicle control has the smallest gap, as fleet ownership cannot affect a one-vehicle decision. Once multiple vehicles must coordinate, COAST improves on the best neural competitor by 5.5\%, 4.7\%, 7.2\%, and 6.1\% from \texttt{n50m3} through \texttt{n400m20}. The gap is not monotone in scale, but it persists after every joint customer--fleet increase; this is consistent with a coordination benefit rather than a simple route-length effect.

\begin{table*}[t]
\centering
\begin{minipage}[t]{0.34\textwidth}
\centering
{\small
\setlength{\tabcolsep}{3.2pt}
\begin{tabular}{lrrr}
\toprule
Scope & $n$ & $\chi^2_F$ & $p_F$ \\
\midrule
All 8 methods & 176 & 1081.14 & $3.52\mathrm{e}{-229}$ \\
Neural subset & 176 & 244.03 & $1.28\mathrm{e}{-52}$ \\
R (all 8) & 63 & 397.49 & $8.23\mathrm{e}{-82}$ \\
C (all 8) & 57 & 336.30 & $1.05\mathrm{e}{-68}$ \\
RC (all 8) & 56 & 352.19 & $4.19\mathrm{e}{-72}$ \\
\bottomrule
\end{tabular}
}\\[1pt]
\footnotesize (a) Friedman tests.
\end{minipage}\hfill
\begin{minipage}[t]{0.62\textwidth}
\centering
{\small
\setlength{\tabcolsep}{3.0pt}
\begin{tabular}{lrrrrr}
\toprule
Baseline & W/L & $\tilde{\Delta}$ & $T_{\min}$ & $p_{\mathrm{raw}}$ & $p_{\mathrm{Holm}}$ \\
\midrule
MARDAM  & 157/19 & 23.4\% & 522 & $7.03\mathrm{e}{-27}$ & $7.03\mathrm{e}{-27}$ \\
AM      & 166/10 & 14.4\% & 206 & $4.03\mathrm{e}{-29}$ & $1.21\mathrm{e}{-28}$ \\
PolyNet & 161/15 & 13.0\% & 408 & $1.12\mathrm{e}{-27}$ & $2.24\mathrm{e}{-27}$ \\
GPHH    & 172/4  & 27.6\% & 27  & $1.97\mathrm{e}{-30}$ & $8.68\mathrm{e}{-30}$ \\
C+C     & 176/0  & 62.7\% & 0   & $1.24\mathrm{e}{-30}$ & $8.68\mathrm{e}{-30}$ \\
C+W     & 176/0  & 62.5\% & 0   & $1.24\mathrm{e}{-30}$ & $8.68\mathrm{e}{-30}$ \\
WIQ+C   & 176/0  & 68.1\% & 0   & $1.24\mathrm{e}{-30}$ & $8.68\mathrm{e}{-30}$ \\
\bottomrule
\end{tabular}
}\\[1pt]
\footnotesize (b) Two-sided Wilcoxon tests.
\end{minipage}
\caption{RQ1 statistical validation on the 176 matched Solomon instances. In (a), ``all 8'' comprises COAST, MARDAM, AM, PolyNet, GPHH, C+C, C+W, and WIQ+C; the neural subset is a like-for-like complement. In (b), W/L is COAST's win/loss count, $\tilde{\Delta}$ is its median percentage cost reduction, and $T_{\min}$ is the smaller of the positive and negative Wilcoxon signed-rank sums. Holm adjustment is applied to the seven prespecified post hoc contrasts shown here.}
\label{tab:rq1_statistics}
\end{table*}

\paragraph{Statistical validation.} Table~\ref{tab:rq1_statistics} reports the inferential evidence after the two benchmark views. The all-method Friedman test rejects equal performance ($\chi^2_F=1081.14$, $p=3.52\mathrm{e}{-229}$); the neural-only family and the R, C, and RC subsets yield the same conclusion. The post hoc panel covers every reported Solomon comparator, including GPHH and all three dispatching rules. Each COAST contrast remains significant after the seven-way Holm correction, including all three neural baselines (MARDAM: $p_{\mathrm{Holm}}=7.03\mathrm{e}{-27}$; AM: $1.21\mathrm{e}{-28}$; PolyNet: $2.24\mathrm{e}{-27}$). These tests establish consistency across matched instances; they do not substitute for multi-seed training validation.

\subsection{RQ2: How do dynamism and time windows change the advantage?}

\paragraph{Question.} Does COAST retain an advantage as customer arrivals become more dynamic and time windows become tighter?

We evaluate a full-factorial \texttt{n50m3} grid with four degrees of dynamism (DoD), four time-window (TW) ratios, and 32 matched instances per cell (512 total). This is a \emph{within-support sensitivity analysis}: it varies factors represented in data construction rather than claiming an out-of-distribution result. Figure~\ref{fig:rq2_components} reports all four neural methods on the complete grid with one shared colour scale, so the cellwise comparison is not restricted to a single baseline.

\begin{figure*}[t]
\centering
\includegraphics[width=\textwidth]{figures/fig_rq2_four_heatmaps.pdf}
\caption{Cellwise DoD--TW comparison of all four neural methods arranged horizontally. Each panel displays the mean normalized cost of one method over 32 matched instances per cell; lower is better. All panels share a common colour scale, enabling direct visual comparison across methods. COAST is lowest in every cell, while MARDAM is the closest comparator throughout this grid.}
\label{fig:rq2_components}
\end{figure*}

Figure~\ref{fig:rq2_components} shows a consistent ordering across the complete grid: COAST has the lowest cost in all 16 cells, MARDAM is next-lowest throughout, and both AM and PolyNet remain higher. MARDAM is therefore the most stringent cellwise comparator. At DoD $=0.10$, COAST's reduction over MARDAM is 0.5--1.5\% across TW ratios; at DoD $=0.75$, it is 2.3--6.2\%. The largest reduction occurs at DoD $=0.75$, TW $=0.25$, and remains 5.6\% and 5.1\% at TW $=0.50$ and $0.75$, respectively. Across the full grid, COAST's reduction ranges from 5.4--13.0\% versus AM and 4.4--13.9\% versus PolyNet. The high-dynamism pattern is therefore not an artifact of selecting one baseline; it is consistent with the greater cost of early decisions when more requests remain to arrive.

The interaction is not monotonic in both factors. At the widest TW ratio ($1.00$), COAST still leads every comparator, but its high-DoD reduction over MARDAM falls to 2.3\%. Thus, the evidence does not support a claim that lookahead or ownership helps uniformly in every condition. Instead, within this grid, the largest COAST margins occur under high dynamism when time windows are not widest. This pattern is consistent with, but does not itself prove, the proposed separation of current vehicle--customer compatibility from fleet coordination and future impact.

The marginal means reinforce the cellwise result. All methods become costlier as TW increases, yet COAST remains nearly flat from DoD $=0.10$ to $0.50$ (43.64 to 43.61) before rising to 44.45 at $0.75$; MARDAM increases monotonically from 44.07 to 46.68. At TW $=1.00$, COAST reaches 45.90, compared with 46.62 for MARDAM and approximately 50.7 for AM and PolyNet. Relative to MARDAM, COAST therefore has a flatter response to arrival dynamism while retaining an advantage throughout the time-window sweep.

\subsection{RQ3: Which components matter, and when?}

\paragraph{Question.} Which decision signals account for COAST's gains, and do their effects become clearer in the regimes they were designed to address?

\begin{table}[t]
\centering
{\small
\setlength{\tabcolsep}{2.4pt}
\begin{tabular}{lrrrr}
\toprule
Control & \texttt{n50m3} & \texttt{n100m5} & \texttt{n200m10} & \texttt{n400m20} \\
\midrule
B5: fixed fusion & +3.6\% & +7.2\% & \textbf{+42.1\%} & \textbf{+44.9\%} \\
B0: base attention & +3.1\% & +2.6\% & +2.3\% & +2.1\% \\
B1: memory-only & +3.4\% & +3.0\% & +2.1\% & +1.5\% \\
B3: lookahead-only & +2.3\% & +2.1\% & +1.3\% & +0.5\% \\
\bottomrule
\end{tabular}
}
\caption{Fusion and structural controls: percentage increase in mean cost over COAST. B5 preserves all signals but fixes their fusion to an equal-weight sum. \texttt{n20m1} is omitted because it has no fleet coordination.}
\label{tab:rq3_controls}
\end{table}

COAST has lower mean cost in all 28 variant--scale comparisons. Figure~\ref{fig:rq3_components}(a) contains only the three direct removals: each line is the mean-cost increase after removing one named signal. Removing ownership increases cost by 1.7--2.9\% at every scale, removing edge features by 2.5--3.4\%, and removing lookahead by 0.3--1.9\%. The direct-removal view shows a consistent scale effect for ownership and edge features; the smaller lookahead effect motivates the targeted regime analysis below rather than a universal claim. B5 is excluded from this panel because it changes the fusion rule rather than removing a signal.

Table~\ref{tab:rq3_controls} addresses a separate question: whether the signals must be fused adaptively. B5 retains the three signal sources but replaces learned fusion with an equal-weight sum. Its cost increase grows from 3.6\% at \texttt{n50m3} to 7.2\%, 42.1\%, and 44.9\% at the three larger scales. Because B5 changes fusion rather than removing a source of information, it tests whether the signals need to be combined adaptively. B0, B1, and B3 show that simpler score compositions are also consistently worse; because each changes several connections at once, they are structural controls rather than source-specific causal tests.

\begin{figure*}[t]
\centering
\begin{minipage}[t]{0.49\textwidth}
\centering
\includegraphics[width=\linewidth]{figures/fig_rq3_scale_column.pdf}\\[-2pt]
\footnotesize (a) Direct signal removals across scale.
\end{minipage}\hfill
\begin{minipage}[t]{0.49\textwidth}
\centering
\includegraphics[width=\linewidth]{figures/fig_rq3_stress_column.pdf}\\[-2pt]
\footnotesize (b) Effects in targeted operating regimes.
\end{minipage}
\caption{Complementary ablation views. In (a), each line removes one signal and reports the mean-cost increase relative to COAST. In (b), each cell reports the percentage mean-cost change over 500 matched evaluation instances after changing the design in its row: Base is in-distribution evaluation; Burst, Scale, Fleet, Sparse, and Tight respectively alter arrival dynamics, problem scale, fleet size, spatial structure, and time-window tightness. Positive (warm) values favor retaining the original design. The rule separates direct removals from the fixed-fusion B5 control.}
\label{fig:rq3_components}
\end{figure*}

Figure~\ref{fig:rq3_components}(b) asks \emph{when} the designs matter. In its first three rows, a positive cell means that removing the named signal increases cost; in the B5 row, it means that replacing adaptive fusion with fixed fusion increases cost. Removing ownership increases cost in every targeted regime (2.3--3.5\%), consistent with its coordination role. Lookahead has its largest effects under burst arrivals (+5.3\%) and tight windows (+5.2\%), the regimes in which an assignment has greater downstream consequences. Edge features have their largest effects under tight windows (+4.0\%) and sparse spatial structure (+4.7\%), where travel and temporal-feasibility cues are more discriminative.

Negative and near-neutral cells are reported rather than masked by an average. Removing lookahead improves the sparse-spatial regime by 2.2\%, and fixed-fusion B5 is near-neutral or favorable in some targeted regimes despite its large degradation on the scale benchmark. These patterns concern distinct evaluation axes: Table~\ref{tab:rq3_controls} varies the standard PYTH scale, whereas Figure~\ref{fig:rq3_components}(b) perturbs operating characteristics. Together, they support a more precise conclusion: ownership, anticipation, and edge-level feasibility are complementary context-sensitive signals, while adaptive fusion is critical for the multi-vehicle scale transfer tested here. They do not establish that each signal helps in every possible operating condition.

% Let the final wide diagnostic float to the next page if needed; forcing it
% here would create a largely empty float-only page before the limitations.

\subsection{Limitations}

The internal comparisons use one trained seed per configuration, so paired tests establish consistency over test instances rather than training-seed robustness. The lookahead head is optimized only through policy gradient, and its context-dependent stress profile motivates multi-seed validation and stronger training signals. Finally, normalized cost is the primary metric; a route-verified decomposition into distance, late service, and skipped requests is deferred to the artifact.

% Legacy experiment draft retained below for result-audit traceability; it is
% intentionally excluded from the submission in favor of the RQ-based section.
\iffalse

We evaluate COAST on two complementary benchmark families to assess routing quality, dynamic robustness, and the contribution of each architectural component. The primary metric is normalized routing cost under greedy decoding; sampling-based decoding is reported where available as a supplementary reference. Unless otherwise stated, lower cost is better.

\subsection{Experimental setup}

Our experimental study is organized around three complementary questions. First, we evaluate whether COAST improves solution quality relative to neural and classical baselines on established DVRPTW benchmark families. Second, we examine robustness under varying degrees of dynamism and time-window tightness to assess whether the proposed decomposition remains effective beyond a single operating regime. Third, we perform targeted ablations to isolate the contributions of ownership, lookahead, edge-aware scoring, and multi-signal fusion. The remainder of this subsection therefore introduces the evaluation metric, benchmark instances, baselines, and training configuration that support these analyses.

\paragraph{Evaluation metric}
For each test instance, we report the total normalized routing cost of the decoded solution. Let $d_t$ denote the travel distance incurred by the acting vehicle at decision step $t$, and let $s_t^{\mathrm{start}}$ be the service start time at the selected customer. If the selected customer has closing time $b_t$, then the stepwise lateness is defined as
\[
\ell_t = \max(0, s_t^{\mathrm{start}} - b_t).
\]
The cost of one instance is
\[
C
=
\sum_t d_t
\;+\;
\lambda_{\mathrm{late}} \sum_t \ell_t
\;+\;
\lambda_{\mathrm{skip}} N_{\mathrm{skip}},
\]
where $N_{\mathrm{skip}}$ is the number of customer requests that remain unserved when the episode terminates. Thus, the metric jointly captures route efficiency, time-window violation, and service completeness. Throughout the experiments, we use $\lambda_{\mathrm{late}} = 1$ and $\lambda_{\mathrm{skip}} = 2$.

The term \emph{normalized} refers to the representation of the benchmark instances prior to inference. Customer coordinates are rescaled to $[0,1]$, demands are normalized by vehicle capacity, and all temporal quantities, including ready times, closing times, service durations, and appearance times, are normalized by the planning horizon; vehicle speed is rescaled consistently. As a result, the reported objective values are dimensionless normalized costs rather than raw physical distances or times.

\paragraph{Benchmark instances}
We employ two groups of benchmark instances. The first group comprises 176 instances derived from Solomon-style topologies, organized by size and spatial distribution: 56 instances with 100 customers, 60 with 200 customers, and 60 with 400 customers. Each size level includes three topology types --- Random (R), Cluster (C), and a mixed Random-Cluster (RC) --- which impose different routing demands. Random instances test spatial dispersion, Cluster instances test local packing efficiency, and RC instances combine both characteristics. We will call this group the solomon-style benchmark for brevity, and refer to the individual datasets as \texttt{h100}, \texttt{h200}, and \texttt{h400}.

The second group consists of five synthetically generated datasets at increasing scales: \texttt{n20m1} (20 customers, 1 vehicle), \texttt{n50m3} (50 customers, 3 vehicles), \texttt{n100m5}, \texttt{n200m10}, and \texttt{n400m20}. These datasets scale both the number of customers and the fleet size simultaneously, enabling evaluation of coordination requirements under growing problem complexity. Each dataset contains 10,240 instances generated from the same distribution used during training. Each dataset is created by combining four degrees of dynamism (DoD = 0.10, 0.25, 0.50, 0.75) with four time-window ratios (0.25, 0.50, 0.75, 1.00). This group will be referred to as the PYTH benchmark, and the individual datasets will be referred to by their \texttt{nXmY} notation.

\paragraph{Baselines}
We compare COAST against several established methods. \textbf{MARDAM} is a multi-head attention-based routing policy that serves as a direct neural baseline with comparable capacity. \textbf{Attention Model (AM)} \cite{kool2018attention} and \textbf{PolyNet} \cite{hottung2025polynet} are Transformer-based routing architectures originally developed for static VRP, adapted here for the dynamic setting. We additionally include classical hyper-heuristic baselines: \textbf{GPHH} (genetic programming hyper-heuristic) \cite{TAM2026114234}, \textbf{C+C} (Clark-Wright savings with competence), \textbf{C+W} (Clark-Wright with waiting), and \textbf{WIQ+C} (work-in-queue plus competence). All neural baselines use the same training data distribution.

\paragraph{Training configuration}
Model-specific architectural settings and parameter counts are summarized in Table~\ref{tab:model_config}. The key point for interpretation is that COAST and all internal ablations share the same backbone, so the ablation study isolates the contribution of the proposed modules rather than capacity differences. By contrast, MARDAM, AM, and PolyNet should be viewed as external neural baselines with comparable or stronger capacity, making improvements over them conservative rather than attributable to an oversized backbone.

Training uses 500 epochs (1,000 iterations each) with a batch size of 512, so the number of training instances is 512,000. The learning rate is $10^{-4}$ via the Adam optimizer. The policy gradient objective employs a learned critic baseline with advantage normalization and entropy regularization ($\beta = 0.01$). All models are trained on DVRPTW instances with 50 customers and 3 vehicles; time windows are drawn uniformly from $[30, 91]$, demands from $[5, 41]$, and the degree of dynamism is sampled uniformly from $\{0.10, 0.25, 0.50, 0.75\}$ per instance, exposing all policies to low-, medium-, and high-dynamism regimes during training rather than tuning them to a single operating point.

\begin{table}[t]
\centering
{\small
\setlength{\tabcolsep}{1.5pt}
\begin{tabular}{lccccc}
\toprule
Model & Layers & Heads & FF & Dim & Params \\
\midrule
\shortstack[l]{COAST + internal\\ablations} & 2 & 4 & 256 & 128 & 0.649M  \\
MARDAM & 3 & 8 & 512 & 128 & 0.727M  \\
AM & 5 & 8 & 256 & 128 & 0.794M  \\
PolyNet & 5 & 8 & 256 & 128 & 0.830M  \\
\bottomrule
\end{tabular}
}
\caption{Neural model configurations used in the experiments. The first row lists COAST and its internal ablations, which share the same backbone architecture. The subsequent rows list external neural baselines with comparable or stronger capacity.}
\label{tab:model_config}
\end{table}



\subsection{Main results on Solomon-style datasets}

Table~\ref{tab:solomon_benchmark} reports the Solomon-style benchmark results grouped by instance size and topology. Instead of using only size-level averages, we separately report Random (R), Cluster (C), and mixed Random-Cluster (RC) instances because these topologies impose different routing challenges. Random instances emphasize global assignment and long-range coordination, Cluster instances favor compact intra-cluster routing, while RC instances require both local exploitation and global coordination.

At the smallest scale, \texttt{h100}, COAST(g) achieves the best mean cost on the Random and RC subsets, with costs of 51.33 and 65.69, respectively. The only exception is the Cluster subset, where MARDAM(g) slightly outperforms COAST(g), obtaining 65.08 compared with 67.65. This suggests that, for small and spatially concentrated instances, the benefit of the proposed decomposition is less pronounced because the routing structure is already relatively compact. Nevertheless, after aggregating all \texttt{h100} topologies, COAST(g) still achieves the best overall mean cost of 60.39.

As the instance size increases, the advantage of COAST becomes more consistent. On \texttt{h200}, COAST(g) obtains the lowest mean cost across all three topology groups, with an aggregate cost of 108.70 compared with 149.27 for MARDAM(g), 124.88 for AM(g), and 129.44 for PolyNet(g). The same trend is further amplified on \texttt{h400}, where COAST(g) again ranks first on Random, Cluster, and RC instances, leading to the best aggregate cost of 208.20. These results indicate that COAST scales more effectively as the number of customers grows, where route decomposition and coordination become increasingly important.

The topology-level results further show that the gains are not concentrated in a single spatial regime. At \texttt{h200}, MARDAM(g) is 37.4\%, 38.1\%, and 36.4\% higher than COAST(g) on the Random, Cluster, and RC groups, respectively, when measured relative to COAST(g). At \texttt{h400}, these gaps increase to 52.5\%, 46.8\%, and 58.8\%. The largest gap appears on the \texttt{h400}-RC subset, suggesting that COAST is particularly effective when both clustered local structure and dispersed global structure must be handled simultaneously.

The instance-level comparison in Table~\ref{tab:solomon_win_counts} confirms that the aggregate improvements are not caused by a small number of outlier instances. Across all 176 Solomon-style instances, COAST(g) outperforms MARDAM(g) on 157 instances, AM(g) on 166 instances, PolyNet(g) on 161 instances, and GPHH on 172 instances. COAST(g) also wins against C+C, C+W, and WIQ+C on all evaluated instances. The win rates become especially strong at larger scales: on both \texttt{h200} and \texttt{h400}, COAST(g) wins against MARDAM(g) on 100.0\% of the instances.

Table~\ref{tab:solomon_improvement} reports the average percentage improvement of COAST(g) over each baseline. Over all instances, COAST(g) improves upon MARDAM(g), AM(g), PolyNet(g), and GPHH by 22.4\%, 13.8\%, 13.9\%, and 33.3\%, respectively. The improvement over MARDAM(g) increases from 2.0\% on \texttt{h100} to 28.6\% on \texttt{h200} and 35.3\% on \texttt{h400}, further supporting the scalability of COAST. Similar trends are observed against the classical and handcrafted baselines, where the improvements exceed 60\% on average for C+C, C+W, and WIQ+C.

Comparing decoding strategies in Table~\ref{tab:solomon_benchmark}, COAST achieves strong results with both sampling and greedy decoding. The difference between COAST(s) and COAST(g) is generally small, and greedy decoding is often competitive or better in the aggregate rows. This indicates that COAST does not rely heavily on stochastic sampling to obtain high-quality solutions, which is desirable for efficient inference.

Overall, the results in Tables~\ref{tab:solomon_benchmark}, \ref{tab:solomon_win_counts}, and \ref{tab:solomon_improvement} show that COAST is competitive at small scale and becomes increasingly advantageous on larger Solomon-style instances. Except for the \texttt{h100} Cluster subset, COAST consistently delivers the best greedy performance across topology groups. The high instance-level win rates and substantial percentage improvements further demonstrate that the proposed method provides robust and scalable gains across different spatial distributions.


\begin{table*}[t]
\centering
{\small
\setlength{\tabcolsep}{1.5pt}
\begin{tabular}{llrrrrrrrrrrrr}
\toprule
& & \multicolumn{2}{c}{COAST} & \multicolumn{2}{c}{MARDAM} & \multicolumn{2}{c}{AM} & \multicolumn{2}{c}{PolyNet} & & & & \\
Dataset & Topology & s & g & s & g & g & s & g & s & GPHH & C+C & C+W & WIQ+C \\
\midrule
\texttt{h100} & Random & 52.73 & \textbf{51.33} & 56.53 & 54.89 & 61.22 & 60.28 & 59.57 & 98.89 & 66.39 & 144.02 & 144.40 & 178.61 \\
\texttt{h100} & Cluster & 66.68 & 67.65 & 66.11 & \textbf{65.08} & 73.01 & 72.56 & 70.36 & 100.38 & 80.63 & 130.71 & 118.12 & 175.02 \\
\texttt{h100} & RC & 66.57 & \textbf{65.69} & 68.90 & 67.42 & 73.07 & 71.80 & 70.31 & 112.22 & 80.11 & 152.99 & 141.10 & 181.66 \\
\texttt{h100} & All & 60.92 & \textbf{60.39} & 62.97 & 61.56 & 68.18 & 67.30 & 65.91 & 103.15 & 74.63 & 142.54 & 135.48 & 178.39 \\
\midrule
\texttt{h200} & Random & 106.03 & \textbf{103.29} & 145.03 & 141.90 & 122.33 & 121.73 & 120.81 & 169.48 & 188.74 & 270.39 & 262.94 & 346.91 \\
\texttt{h200} & Cluster & 118.06 & \textbf{117.61} & 163.22 & 162.43 & 126.64 & 127.77 & 137.04 & 171.67 & 191.74 & 269.58 & 248.67 & 361.24 \\
\texttt{h200} & RC & \textbf{104.97} & 105.19 & 146.90 & 143.48 & 125.66 & 126.27 & 130.48 & 183.64 & 181.00 & 282.86 & 283.12 & 350.99 \\
\texttt{h200} & All & 109.69 & \textbf{108.70} & 151.72 & 149.27 & 124.88 & 125.26 & 129.44 & 174.93 & 187.16 & 274.28 & 264.91 & 353.05 \\
\midrule
\texttt{h400} & Random & 208.86 & \textbf{200.71} & 317.68 & 306.15 & 237.63 & 238.96 & 232.22 & 349.95 & 356.53 & 613.33 & 600.64 & 716.09 \\
\texttt{h400} & Cluster & 225.12 & \textbf{220.65} & 339.35 & 323.84 & 248.35 & 251.41 & 258.78 & 361.37 & 374.85 & 630.06 & 616.13 & 752.28 \\
\texttt{h400} & RC & 208.20 & \textbf{203.26} & 341.53 & 322.73 & 250.32 & 253.20 & 244.98 & 405.95 & 352.72 & 634.84 & 636.76 & 724.18 \\
\texttt{h400} & All & 214.06 & \textbf{208.20} & 332.86 & 317.57 & 245.43 & 247.86 & 245.33 & 372.42 & 361.37 & 626.07 & 617.85 & 730.85 \\
\midrule
All & All & 129.75 & \textbf{127.25} & 185.23 & 178.74 & 147.94 & 148.61 & 148.73 & 219.42 & 210.74 & 352.29 & 344.05 & 426.27 \\
\bottomrule
\end{tabular}
}
\caption{Solomon-style benchmark results grouped by dataset size and topology. Values are mean normalized costs over the evaluated scenarios in each row. The \texttt{All} row aggregates all topologies within the same dataset size, and the final row aggregates all 176 Solomon-style instances.}
\label{tab:solomon_benchmark}
\end{table*}


\begin{figure*}[t]
\centering
\includegraphics[width=0.85\textwidth]{figures/fig_csv_benchmark.pdf}
\caption{Solomon-style benchmark comparison grouped by instance size.}
\label{fig:solomon_benchmark}
\end{figure*}


\begin{table*}[t]
\centering
{\small
\setlength{\tabcolsep}{1.5pt}
\begin{tabular}{lrrrrrrr}
\toprule
Dataset & MARDAM(g) & AM(g) & PolyNet(g) & GPHH & C+C & C+W & WIQ+C \\
\midrule
\texttt{h100} & \shortstack{37/56\\(66.1\%)} & \shortstack{48/56\\(85.7\%)} & \shortstack{43/56\\(76.8\%)} & \shortstack{53/56\\(94.6\%)} & \shortstack{56/56\\(100.0\%)} & \shortstack{56/56\\(100.0\%)} & \shortstack{56/56\\(100.0\%)} \\
\texttt{h200} & \shortstack{60/60\\(100.0\%)} & \shortstack{58/60\\(96.7\%)} & \shortstack{59/60\\(98.3\%)} & \shortstack{59/60\\(98.3\%)} & \shortstack{60/60\\(100.0\%)} & \shortstack{60/60\\(100.0\%)} & \shortstack{60/60\\(100.0\%)} \\
\texttt{h400} & \shortstack{60/60\\(100.0\%)} & \shortstack{60/60\\(100.0\%)} & \shortstack{59/60\\(98.3\%)} & \shortstack{60/60\\(100.0\%)} & \shortstack{60/60\\(100.0\%)} & \shortstack{60/60\\(100.0\%)} & \shortstack{60/60\\(100.0\%)} \\
All & \shortstack{157/176\\(89.2\%)} & \shortstack{166/176\\(94.3\%)} & \shortstack{161/176\\(91.5\%)} & \shortstack{172/176\\(97.7\%)} & \shortstack{176/176\\(100.0\%)} & \shortstack{176/176\\(100.0\%)} & \shortstack{176/176\\(100.0\%)} \\
\bottomrule
\end{tabular}
}
\caption{Instance-level win counts and win rates of COAST(g) on the Solomon-style benchmark. Each entry has the format wins/instances (win rate), where a win indicates that COAST(g) obtains a lower normalized cost than the corresponding baseline.}
\label{tab:solomon_win_counts}
\end{table*}


\begin{table*}[t]
\centering
{\small
\setlength{\tabcolsep}{2pt}
\begin{tabular}{lrrrrrrr}
\toprule
Dataset & MARDAM(g) & AM(g) & PolyNet(g) & GPHH & C+C & C+W & WIQ+C \\
\midrule
\texttt{h100} & 2.0 & 11.0 & 8.0 & 18.3 & 56.9 & 53.1 & 66.1 \\
\texttt{h200} & 28.6 & 14.0 & 17.3 & 40.0 & 61.1 & 59.5 & 69.4 \\
\texttt{h400} & 35.3 & 16.3 & 16.2 & 40.6 & 67.4 & 66.9 & 71.7 \\
All & 22.4 & 13.8 & 13.9 & 33.3 & 61.9 & 60.0 & 69.1 \\
\bottomrule
\end{tabular}
}
\caption{Average percentage improvement of COAST(g) over each baseline on the Solomon-style benchmark. Values are computed per instance as $(\mathrm{baseline}-\mathrm{COAST})/\mathrm{baseline}\times 100$ and then averaged within each group.}
\label{tab:solomon_improvement}
\end{table*}


\subsection{Main results on PYTH datasets}

Tables~\ref{tab:pyth_benchmark} and \ref{tab:pyth_benchmark_baselines} report the results on the five synthetically generated PYTH datasets, while Figure~\ref{fig:pyth_scale} summarizes the scaling trend under greedy decoding. Complementing the Solomon-style experiments, which evaluate robustness across spatial topologies, the PYTH benchmark focuses on scalability by increasing both the number of customers and the fleet size. This setting is therefore useful for examining whether the advantage of COAST persists when route construction and inter-vehicle coordination become increasingly difficult.

Across all PYTH scales, COAST(g) achieves the lowest normalized cost among the neural baselines. On the smallest single-vehicle setting, \texttt{n20m1}, COAST(g) obtains 24.78, slightly improving over MARDAM(g), AM(g), and PolyNet(g). Since this setting contains only one vehicle, it provides limited evidence for fleet-level coordination and mainly tests basic route construction. The comparison becomes more informative from \texttt{n50m3} onward, where multiple vehicles must be coordinated. In these multi-vehicle settings, COAST(g) consistently outperforms the baselines, achieving costs of 43.92, 94.85, 182.38, and 355.99 on \texttt{n50m3}, \texttt{n100m5}, \texttt{n200m10}, and \texttt{n400m20}, respectively.

The performance gap generally widens as the problem scale increases. Relative to COAST(g), MARDAM(g) is 5.4\%, 5.8\%, 4.9\%, 11.4\%, and 20.9\% higher from \texttt{n20m1} to \texttt{n400m20}, respectively. A similar advantage is observed over AM(g) and PolyNet(g), whose costs are consistently higher across all scales; on \texttt{n50m3}, for example, AM(g) and PolyNet(g) remain 12.9\% and 13.7\% above COAST(g). This trend supports the same scalability conclusion observed in the Solomon-style benchmark: COAST remains competitive on smaller instances, but its advantage becomes clearer when the routing problem requires stronger coordination over a larger set of customers and vehicles.

Comparing decoding strategies, COAST(g) is competitive with or better than COAST(s) on four of the five scales and is already stronger from \texttt{n50m3} onward. Sampling is only marginally better on \texttt{n20m1}, where the problem reduces to a compact single-route regime with little coordination burden. This again suggests that COAST does not depend heavily on stochastic sampling to obtain high-quality solutions, making greedy decoding an efficient and reliable inference option for the larger multi-vehicle settings that matter most in this study.

Overall, the PYTH results reinforce the main finding from the Solomon-style experiments. COAST provides stable gains across different benchmark settings, and its advantage becomes more pronounced as routing complexity increases. While the Solomon-style results demonstrate robustness across spatial distributions, the PYTH results show that COAST also scales effectively when both customer count and fleet size grow.

\begin{table*}[t]
\centering
{\small
\setlength{\tabcolsep}{2pt}
\begin{tabular}{lrrrr}
\toprule
& \multicolumn{2}{c}{COAST} & \multicolumn{2}{c}{MARDAM} \\
Dataset & s & g & s & g \\
\midrule
\texttt{n20m1} & $\textbf{24.72}\pm2.40$ & $24.78\pm2.38$ & $26.10\pm3.09$ & $26.13\pm3.16$ \\
\texttt{n50m3} & $44.10\pm4.09$ & $\textbf{43.92}\pm4.09$ & $46.85\pm4.71$ & $46.48\pm4.76$ \\
\texttt{n100m5} & $95.49\pm5.88$ & $\textbf{94.85}\pm5.89$ & $101.06\pm6.70$ & $99.50\pm6.56$ \\
\texttt{n200m10} & $184.70\pm8.95$ & $\textbf{182.38}\pm8.90$ & $209.02\pm14.70$ & $203.24\pm14.43$ \\
\texttt{n400m20} & $363.35\pm13.82$ & $\textbf{355.99}\pm13.43$ & $443.86\pm34.78$ & $430.48\pm35.05$ \\
\bottomrule
\end{tabular}
}
\caption{PYTH benchmark results for COAST and MARDAM. Values are mean normalized cost $\pm$ standard deviation.}
\label{tab:pyth_benchmark}
\end{table*}

\begin{table*}[t]
\centering
{\small
\setlength{\tabcolsep}{2pt}
\begin{tabular}{lrrrr}
\toprule
& \multicolumn{2}{c}{AM} & \multicolumn{2}{c}{PolyNet} \\
Dataset & s & g & s & g \\
\midrule
\texttt{n20m1} & $25.74\pm3.13$ & $25.40\pm2.75$ & $31.84\pm4.08$ & $25.58\pm2.94$ \\
\texttt{n50m3} & $50.28\pm5.49$ & $49.58\pm5.00$ & $68.57\pm8.22$ & $49.93\pm5.31$ \\
\texttt{n100m5} & $106.18\pm8.06$ & $103.85\pm7.14$ & $144.53\pm12.87$ & $105.13\pm8.42$ \\
\texttt{n200m10} & $201.54\pm13.02$ & $196.53\pm10.92$ & $293.54\pm24.83$ & $200.89\pm14.36$ \\
\texttt{n400m20} & $389.99\pm21.35$ & $378.98\pm17.10$ & $558.21\pm41.05$ & $389.72\pm26.42$ \\
\bottomrule
\end{tabular}
}
\caption{PYTH benchmark results for AM and PolyNet. Values are mean normalized cost $\pm$ standard deviation.}
\label{tab:pyth_benchmark_baselines}
\end{table*}

\begin{figure}[t]
\centering
\includegraphics[width=\linewidth]{figures/fig_pyth_scale.pdf}
\caption{Scalability on PYTH datasets under greedy decoding.}
\label{fig:pyth_scale}
\end{figure}



\subsection{Effect of dynamism}

To isolate the effect of online uncertainty, we evaluate a dedicated dynamic benchmark grid built from the same \texttt{n50m3} regime used in training: 50 customers, 3 vehicles, horizon 480, and greedy decoding. The grid contains 512 instances organized as 16 cells, obtained by combining four degrees of dynamism (DoD = 0.10, 0.25, 0.50, 0.75) with four time-window ratios (0.25, 0.50, 0.75, 1.00), with 32 instances per cell. Table~\ref{tab:dynamic_summary} reports the aggregated cell averages by DoD and by time-window ratio, while Figure~\ref{fig:dynamic_sensitivity} provides cell-level heatmaps over the full DoD $\times$ TW grid for all four neural methods.

The full grid confirms that COAST is the best-performing method in all 16 DoD--TW cells. Averaged over the entire grid, COAST attains a mean normalized cost of 43.80, compared with 45.02 for MARDAM, 47.15 for PolyNet, and 47.65 for AM. The best COAST cell occurs at DoD = 0.50 and TW = 0.25 with cost 41.21, while the most difficult cell is DoD = 0.75 and TW = 1.00 with cost 46.88. Even in that hardest setting, COAST remains better than MARDAM (47.96), PolyNet (49.44), and AM (49.53), indicating that the advantage is not confined to easy operating points.

As DoD increases from 0.10 to 0.75, COAST remains comparatively stable, with mean cost increasing only from 43.64 to 44.45, a range of 0.81. Over the same sweep, MARDAM increases from 44.07 to 46.68, a much larger range of 2.61, while PolyNet and AM average 47.78 to 46.85 and 48.08 to 47.29, respectively. COAST therefore preserves the best absolute performance throughout the entire uncertainty range, and its advantage over MARDAM widens from 1.0\% at DoD = 0.10 to 5.0\% at DoD = 0.75 when averaged over time-window settings. At the cell level, the gap to MARDAM ranges from only 0.5\% at (DoD = 0.10, TW = 0.50) to 6.7\% at (DoD = 0.75, TW = 0.25), which is consistent with the intended role of the ownership and lookahead signals: when more customers are revealed online, the policy benefits more from explicit coordination and future-aware scoring.

All methods also degrade as the time-window ratio increases. Averaged over DoD, COAST rises from 41.71 at TW = 0.25 to 45.90 at TW = 1.00, while MARDAM rises from 43.51 to 46.62, PolyNet from 44.01 to 50.76, and AM from 44.83 to 50.73. This confirms that tighter windows reduce the feasible action set for every model, but the degradation is consistently milder for COAST than for the stronger neural baselines. The gap to AM, for example, is 7.5\% at TW = 0.25 and remains 10.5\% at TW = 1.00; against PolyNet, the corresponding gaps are 5.5\% and 10.6\%. Taken together, these results show that the proposed architecture is not only more robust to increasing dynamism, but also better at preserving route quality when temporal feasibility becomes more restrictive.

\begin{table}[t]
\centering
{\small
\setlength{\tabcolsep}{1.5pt}
\begin{tabular}{llrrrr}
\toprule
Factor & Level & COAST & MARDAM & PolyNet & AM \\
\midrule
DoD & 0.10 & \textbf{43.64} & 44.07 & 47.78 & 48.08 \\
DoD & 0.25 & \textbf{43.50} & 44.35 & 47.25 & 47.68 \\
DoD & 0.50 & \textbf{43.61} & 44.99 & 46.73 & 47.55 \\
DoD & 0.75 & \textbf{44.45} & 46.68 & 46.85 & 47.29 \\
\midrule
TW & 0.25 & \textbf{41.71} & 43.51 & 44.01 & 44.83 \\
TW & 0.50 & \textbf{43.15} & 44.37 & 45.89 & 46.73 \\
TW & 0.75 & \textbf{44.43} & 45.59 & 47.94 & 48.31 \\
TW & 1.00 & \textbf{45.90} & 46.62 & 50.76 & 50.73 \\
\bottomrule
\end{tabular}
}
\caption{Dynamic benchmark summary aggregated by degree of dynamism and time-window ratio. Values are mean normalized cost over the corresponding four cells.}
\label{tab:dynamic_summary}
\end{table}

\begin{figure}[t]
\centering
\includegraphics[width=\linewidth]{figures/fig_dynamic_sensitivity.pdf}
\caption{Cell-level heatmaps on the dynamic benchmark grid. Each panel shows the mean normalized cost of one neural method for every DoD $\times$ TW configuration; lower values are better (lighter is better).}
\label{fig:dynamic_sensitivity}
\end{figure}

\subsection{Ablation study}

To isolate the contribution of each architectural component, we evaluate several ablated variants on the PYTH datasets (Tables~\ref{tab:ablation_pyth} and \ref{tab:ablation_modules}). The ablations systematically remove or simplify key components: \textbf{B0} removes all structured signals (no memory, ownership, or lookahead), leaving only the base attention mechanism; \textbf{B1} retains memory-based vehicle state but removes ownership and lookahead; \textbf{B3} includes only the lookahead signal without coordination; \textbf{B5} retains all signals but replaces the learned non-linear fusion with a linear combination; \textbf{EdgeOff} disables explicit edge features while keeping all other components; and two clean ablations remove either the ownership head or the lookahead head individually.

The full COAST model demonstrates consistent advantages once the problem involves sufficient customers and vehicles for coordination to matter. From \texttt{n50m3} to \texttt{n400m20}, COAST outperforms all ablated variants. Removing the ownership signal increases normalized cost by 1.7--2.9\% on medium and large datasets, confirming that explicit fleet-level coordination provides measurable benefit. Removing lookahead produces a smaller but consistent increase of 0.3--1.9\%, indicating that anticipation contributes positively but its standalone effect is more modest than coordination.

Edge features also prove consistently beneficial: EdgeOff incurs a 2.5--3.4\% cost increase relative to COAST from \texttt{n50m3} to \texttt{n400m20}. This supports the design choice of exposing constraint-related information --- travel time, waiting time, lateness penalty, temporal slack, feasibility status, and capacity gap --- directly to the scoring function rather than relying solely on latent node embeddings.

The most pronounced ablation effect is observed for B5, which replaces learned non-linear fusion with a linear combination of the three signals. The degradation grows sharply with problem scale: B5 increases cost by 7.2\% on \texttt{n100m5}, 42.1\% on \texttt{n200m10}, and 44.9\% on \texttt{n400m20}. This pattern indicates that the three signals are not simply additive at scale; their interactions must be learned to enable context-dependent trade-offs between immediate compatibility, vehicle ownership, and future routing consequences. On the smallest \texttt{n20m1} configuration, several simplified variants perform competitively, confirming that this regime lacks the multi-vehicle dynamics necessary to distinguish architectural choices.

\begin{table*}[t]
\centering
{\small
\setlength{\tabcolsep}{2pt}
\begin{tabular}{lrrrrr}
\toprule
Dataset & COAST & B0 & B1 & B3 & B5 \\
\midrule
\texttt{n20m1} & $24.78\pm2.38$ & $\textbf{23.51}\pm2.36$ & $23.81\pm2.43$ & $23.45\pm2.43$ & $23.79\pm2.42$ \\
\texttt{n50m3} & $\textbf{43.92}\pm4.09$ & $45.28\pm4.26$ & $45.43\pm4.32$ & $44.91\pm4.25$ & $45.49\pm4.31$ \\
\texttt{n100m5} & $\textbf{94.85}\pm5.89$ & $97.35\pm6.17$ & $97.71\pm6.36$ & $96.86\pm6.15$ & $101.67\pm9.02$ \\
\texttt{n200m10} & $\textbf{182.38}\pm8.90$ & $186.52\pm9.21$ & $186.26\pm9.63$ & $184.75\pm9.13$ & $259.10\pm34.30$ \\
\texttt{n400m20} & $\textbf{355.99}\pm13.43$ & $363.43\pm13.86$ & $361.33\pm14.61$ & $357.81\pm13.77$ & $515.72\pm67.45$ \\
\bottomrule
\end{tabular}
}
\caption{PYTH ablation results for the structured-signal variants. Values are mean normalized cost $\pm$ standard deviation.}
\label{tab:ablation_pyth}
\end{table*}

\begin{table*}[t]
\centering
{\small
\setlength{\tabcolsep}{2pt}
\begin{tabular}{lrrrr}
\toprule
Dataset & COAST & EdgeOff & \shortstack{no\\ownership} & \shortstack{no\\lookahead} \\
\midrule
\texttt{n20m1} & $24.78\pm2.38$ & $24.91\pm2.34$ & $23.37\pm2.46$ & $23.96\pm2.49$ \\
\texttt{n50m3} & $\textbf{43.92}\pm4.09$ & $45.43\pm4.20$ & $45.20\pm4.25$ & $44.74\pm4.22$ \\
\texttt{n100m5} & $\textbf{94.85}\pm5.89$ & $97.25\pm5.99$ & $97.44\pm6.27$ & $96.44\pm6.17$ \\
\texttt{n200m10} & $\textbf{182.38}\pm8.90$ & $186.94\pm8.90$ & $186.22\pm9.26$ & $184.33\pm9.43$ \\
\texttt{n400m20} & $\textbf{355.99}\pm13.43$ & $365.43\pm13.16$ & $362.11\pm13.56$ & $357.18\pm14.44$ \\
\bottomrule
\end{tabular}
}
\caption{PYTH ablation results for edge features and individual decision heads. Values are mean normalized cost $\pm$ standard deviation.}
\label{tab:ablation_modules}
\end{table*}

\subsection{Hypothesis analysis}

The proposed method is built around four architectural claims: \textbf{(H1)} explicit ownership improves fleet coordination, \textbf{(H2)} candidate-conditioned lookahead reduces myopic dispatching, \textbf{(H3)} edge-aware scoring improves constraint handling, and \textbf{(H4)} non-linear fusion is needed to combine these signals robustly. These hypotheses are not tested by a single benchmark alone. Instead, we validate them using complementary evidence from the PYTH scale experiments, the targeted ablations in Tables~\ref{tab:ablation_pyth} and \ref{tab:ablation_modules}, and the dynamic-grid analysis in Table~\ref{tab:dynamic_summary} and Figure~\ref{fig:dynamic_sensitivity}. This design is important because each hypothesis concerns a different failure mode: coordination across vehicles, anticipation of future consequences, sensitivity to edge-level feasibility, and interaction among multiple decision signals.

Table~\ref{tab:hypotheses} summarizes the mapping from hypothesis to comparison. For H1 and H2, the cleanest tests are direct removals of the ownership and lookahead modules, respectively. For H3, the relevant comparison is COAST versus EdgeOff, especially under the tighter time-window and higher-dynamism cells where feasibility pressure is strongest. For H4, the key comparison is COAST versus B5, which keeps all signals but replaces the learned MLP fusion with a linear combination. Throughout this subsection, we emphasize \texttt{n50m3} and larger instances, because \texttt{n20m1} contains only one vehicle and therefore cannot serve as strong evidence for coordination-oriented claims.

\textbf{H1 (ownership improves coordination)} asks whether the explicit ownership mechanism contributes beyond the shared policy backbone. The evidence is strongest in the clean COAST versus no-ownership comparison: from \texttt{n50m3} onward, removing ownership increases normalized cost by 1.7--2.9\%. This confirms that the per-vehicle coordination memory and learned ownership bias contribute to more efficient fleet allocation once multiple vehicles must compete for dynamically revealed customers.

\textbf{H2 (lookahead reduces myopia)} tests whether scoring candidate actions with a future-impact signal improves decision quality. This hypothesis receives positive but more modest support. Removing lookahead consistently increases cost on medium and large PYTH datasets, but the effect remains smaller (0.3--1.9\%) than those observed for ownership removal or fusion simplification. The directional trend is therefore stable, but the relatively small magnitude suggests that lookahead helps as a refinement signal rather than as the dominant driver of performance.

\textbf{H3 (edge features improve feasibility)} evaluates whether exposing route-level and constraint-level information directly to the scorer is beneficial. The EdgeOff ablation is consistently worse than COAST, with a 2.5--3.4\% increase on the PYTH scale experiments, and the dynamic benchmark further shows that COAST preserves a clearer advantage under tighter time windows and higher dynamism. Taken together, these results support the claim that explicit edge-level features improve the policy's ability to act under feasibility pressure rather than relying only on latent node embeddings.

\textbf{H4 (non-linear fusion improves robustness)} concerns the interaction among the three signals rather than any signal in isolation. This is the strongest hypothesis in the paper. The B5 ablation, which preserves the same signal sources but replaces the learned non-linear fusion with a linear combination, incurs only modest degradation on very small instances but collapses on larger ones (42.1\% increase on \texttt{n200m10} and 44.9\% on \texttt{n400m20}). This indicates that the signals are not merely additive: their relative importance depends on context, and a linear rule cannot express the required trade-offs between immediate compatibility, ownership, and downstream routing consequences.

Overall, the hypothesis tests support the central decomposition strategy, but with different strengths. H4 is the clearest architectural result because the degradation is both large and strongly scale-dependent. H1 and H3 provide consistent medium-strength evidence across the regimes where coordination and feasibility are genuinely challenging. H2 is directionally supported, but the weaker margin suggests that improving the training signal for lookahead remains an important direction for future work.

\begin{table*}[t]
\centering
{\small
\setlength{\tabcolsep}{1pt}
\begin{tabular}{p{0.20\textwidth}p{0.16\textwidth}p{0.38\textwidth}p{0.20\textwidth}}
\toprule
Hypothesis & Comparison & Main evidence & Conclusion \\
\midrule
H1: ownership improves coordination & COAST vs no\_ownership & no\_ownership is 1.7--2.9\% worse on \texttt{n50m3}--\texttt{n400m20} & Supported on medium/large instances \\
H2: lookahead reduces myopia & COAST vs no\_lookahead & no\_lookahead is 0.3--1.9\% worse on \texttt{n50m3}--\texttt{n400m20} & Mildly supported; effect is smaller \\
H3: edge features improve feasibility & COAST vs EdgeOff & EdgeOff is 2.5--3.4\% worse on PYTH scale; tight-TW cost is 48.00 vs 46.13 & Supported, especially under tight constraints \\
H4: nonlinear fusion improves robustness & COAST vs B5 & B5 is 7.2\% worse on \texttt{n100m5}, 42.1\% on \texttt{n200m10}, and 44.9\% on \texttt{n400m20} & Strongly supported at scale \\
\bottomrule
\end{tabular}
}
\caption{Hypothesis evaluation summary. Positive gaps indicate that the ablated variant incurs higher cost than COAST.}
\label{tab:hypotheses}
\end{table*}

Overall, the results demonstrate that COAST is most effective in the regimes targeted by its architecture: medium-to-large dynamic routing problems where vehicle coordination, temporal feasibility, and action-conditioned anticipation interact. The small single-vehicle regime (\texttt{n20m1}) is less favorable to the full model, which is expected given that coordination requirements are minimal in this setting.

\subsection{Limitations}

Several limitations of the current study should be acknowledged. First, the reported results are based on single-seed training for each model configuration. The smaller effect sizes observed for the lookahead and ownership ablations could be sensitive to training variance; multi-seed experiments with statistical significance testing would strengthen confidence in these findings. Second, the lookahead head is trained solely through the policy gradient objective; a supervised calibration target could potentially enhance its effectiveness and provide a clearer signal for validating the anticipation hypothesis. Third, the current evaluation focuses on normalized routing cost as the primary metric; a more detailed decomposition into travel distance, skipped requests, and time-window penalties would provide additional insight into behavioral differences between methods. Finally, classical optimization baselines such as LKH3 and OR-Tools, while informative for static or oracle settings, are not directly applicable to the dynamic sequential decision framework without substantial modification; their inclusion remains an avenue for future work.

\fi

\section{Conclusions and Future works} \label{sec:futurework}
% Avoid stretching the final, partially filled references page.
\raggedbottom
\bibliography{ref}
\end{document}
